\documentclass[a4paper]{article}

% Layout
\usepackage[margin=1in]{geometry}

\usepackage{microtype}
\usepackage{enumitem}
\usepackage[toc,page]{appendix}

% Mathematics
\usepackage{amsmath,amssymb,amsfonts,amsthm,mathtools}
\usepackage{mathrsfs}
\usepackage{physics}
\usepackage{cancel}

% Graphics
\usepackage{xcolor}
\usepackage{tikz-cd}
\usepackage{pgfplots}
\pgfplotsset{compat=1.18}

% Fonts
\usepackage[T1]{fontenc}
% \usepackage{tgschola}
% \usepackage[scaled=0.92]{helvet}

% \usepackage[
%    lite,
%    subscriptcorrection,
%    slantedGreek,
%    nofontinfo
% ]{mtpro2}

% \usepackage{upgreek}
\usepackage{dsfont}

% Paragraph typography
\setlength{\parskip}{0pt}
\setlength{\parindent}{1.25em}

% Utilities
\usepackage{etoolbox}
\usepackage{dynkin-diagrams}

% Warnings
\usepackage[l2tabu,orthodox]{nag}

% Environments
\newtheorem{theorem}{Theorem}[section]
\newtheorem{lemma}[theorem]{Lemma}
\newtheorem{proposition}[theorem]{Proposition}
\newtheorem{corollary}[theorem]{Corollary}
\newtheorem{conjecture}{Conjecture}

\theoremstyle{definition}
\newtheorem{definition}[theorem]{Definition}
\newtheorem{example}[theorem]{Example}

\theoremstyle{remark}
\newtheorem{remark}[theorem]{Remark}
\newtheorem{outline}[theorem]{Outline}

% Dynkin Diagrams
\def\row#1/#2!{\dynkin{#1}{#2}\\}

\newcommand{\tble}[1]{
   \renewcommand*\do[1]{\row##1!}
   \[
      \begin{array}{ll}
         \docsvlist{#1}
      \end{array}
   \]
}

% Extra Fonts
\DeclareMathAlphabet{\mathpzc}{OT1}{pzc}{m}{it}

\newcommand{\goth}[1]{\mathfrak{#1}}
\newcommand{\pzc}[1]{\mathpzc{#1}}

% Commands for FONT characters
\let\SS\S
\let\ag\gg

\ExplSyntaxOn

% Goth Font
\clist_map_inline:nn
{a,b,c,d,e,f,g,h,i,j,k,l,m,n,o,p,q,r,s,t,u,v,w,x,y,z,A,B,C,D,E,F,G,H,I,J,K,L,M,N,O,P,Q,R,S,T,U,V,W,X,Y,Z}
{
   \cs_undefine:c { g#1 }
   \cs_new:cpn { g#1 } { \goth{#1} }
}

% Scr Font
\clist_map_inline:nn
{A,B,C,D,E,F,G,H,I,J,K,L,M,N,O,P,Q,R,S,T,U,V,W,X,Y,Z}
{
   \cs_new:cpn { f#1 } { \mathcal{#1} }
}

% Cal Font
\clist_map_inline:nn
{A,B,C,D,E,F,G,H,I,J,K,L,M,N,O,P,Q,R,S,T,U,V,W,X,Y,Z}
{
   \cs_new:cpn { s#1 } { \mathcal{#1} }
}

% mathbf Font
\clist_map_inline:nn
{A,B,C,D,E,F,G,H,I,J,K,L,M,N,O,P,Q,R,S,T,U,V,W,X,Y,Z}
{
   \cs_undefine:c { #1 }
   \cs_new:cpn { q#1 } { \mathbf{#1} }
}

% BB Font
\clist_map_inline:nn
{A,B,C,D,E,F,G,H,I,J,K,L,M,N,O,P,Q,R,S,T,U,V,W,X,Y,Z}
{
   \cs_undefine:c { #1 }
   \cs_new:cpn { #1 } { \mathbb{#1} }
}

\ExplSyntaxOff

% Categories
\newcommand{\cC}{\mathbf{C}}
\newcommand{\cSet}{\mathbf{Set}}
\newcommand{\cAb}{\mathbf{Ab}}
\newcommand{\cGroup}{\mathbf{Grp}}
\newcommand{\ckSch}[1]{\mathbf{Sch}(#1)}
\newcommand{\ckAlg}[1]{\mathbf{Alg}(#1)}
\newcommand{\cMod}[1]{\mathbf{Mod}(#1)}
\newcommand{\cdh}{h}
\newcommand{\opp}[1]{{#1}^\text{op}}
\newcommand{\et}[1]{{#1}_\text{ét}}
\newcommand{\sep}[1]{{#1}^\text{sep}}
\newcommand{\mgbar}[1]{\overline{\fM}_{#1}}
\DeclareMathOperator{\cSch}{\mathbf{Sch}}

\DeclareMathOperator{\Mor}{Mor}

% Relative stuff
\DeclareMathOperator{\rSpec}{\mathpzc{Spec}}
\DeclareMathOperator{\rHom}{\mathpzc{Hom}}
\DeclareMathOperator{\rEnd}{\mathpzc{End}}
\DeclareMathOperator{\rProj}{\mathpzc{Proj}}
\DeclareMathOperator{\rExt}{\mathpzc{Ext}}

\DeclareMathOperator{\edim}{edim}

% Linear Algebra/Abelian Category/General Notions (?)
\DeclareMathOperator{\rk}{rank}
\DeclareMathOperator{\im}{im}
\DeclareMathOperator{\coker}{coker}
\DeclareMathOperator{\End}{End}
\DeclareMathOperator{\Hom}{Hom}
\DeclareMathOperator{\Aut}{Aut}
\DeclareMathOperator{\len}{length}
\DeclareMathOperator{\height}{height}
\DeclareMathOperator{\codim}{codim}
\DeclareMathOperator{\id}{id}

% AG/Commutative Algebra/Algebra
\DeclareMathOperator{\Hilb}{Hilb}
\DeclareMathOperator{\Def}{Def}
\DeclareMathOperator{\car}{char}
\DeclareMathOperator{\proj}{Proj}
\DeclareMathOperator{\spec}{Spec}
\DeclareMathOperator{\pic}{Pic}
\DeclareMathOperator{\cacl}{CaCl}
\DeclareMathOperator{\cl}{Cl}
\DeclareMathOperator{\supp}{Supp}
\DeclareMathOperator{\Gr}{Gr}
\DeclareMathOperator{\Div}{Div}
\DeclareMathOperator{\Der}{Der}
\DeclareMathOperator{\Num}{Num}
\DeclareMathOperator{\Ann}{Ann}
\DeclareMathOperator{\Syl}{Syl}
\DeclareMathOperator{\Sym}{Sym}
\DeclareMathOperator{\Stab}{Stab}
\DeclareMathOperator{\Perm}{Perm}
\DeclareMathOperator{\Orb}{Orb}
\DeclareMathOperator{\Gal}{Gal}
\DeclareMathOperator{\Supp}{Supp}
\DeclareMathOperator{\trd}{tr.d.}
\DeclareMathOperator{\Frac}{Frac}

\newcommand{\red}[1]{{#1}_\text{red}}
\newcommand{\gsecx}[1]{\Gamma(X, #1)}
\newcommand{\sh}[1]{{#1}^\text{sh}}

% Homological Algebra
\DeclareMathOperator{\Ext}{Ext}
\DeclareMathOperator{\Tor}{Tor}
\DeclareMathOperator{\hd}{hd}
\DeclareMathOperator{\pd}{pd}
\DeclareMathOperator{\cH}{H}
\DeclareMathOperator{\Ti}{T}
\DeclareMathOperator{\dR}{R}
\DeclareMathOperator{\depth}{depth}

% Linear Algebraic Groups
\DeclareMathOperator{\PGL}{PGL}
\DeclareMathOperator{\SO}{SO}
\DeclareMathOperator{\SU}{SU}
\DeclareMathOperator{\Sp}{Sp}
\DeclareMathOperator{\Orth}{O}
\DeclareMathOperator{\Uni}{U}

% Random Stuff
\DeclareMathOperator{\Hasse}{Hasse}
\newcommand{\dlog}{\dd \log}

% Miscellaneous
\author{James Lee}

\allowdisplaybreaks

% Solarized
% \pagecolor[RGB]{0,20,26}
% \color[RGB]{191,191,191}

% Sephia
% \pagecolor[RGB]{249,239,220}

% Grey
% \pagecolor[RGB]{200, 200, 200}

% Black
% \pagecolor[RGB]{24,24,24}
% \color[RGB]{170,180,190}

% Cyber Deep Sea
% \pagecolor[RGB]{10,18,24}
% \color[RGB]{176,204,208}

\newcommand\restr[2]{{\left.\kern-\nulldelimiterspace #1 \vphantom{\big|} \right|_{#2}}}

\let\oldbigoplus\bigoplus

\renewcommand{\bigoplus}{\mathop{\textstyle\oldbigoplus}\displaylimits}


\title{Chapter 5, Section 6}

\begin{document}

\maketitle

\begin{enumerate} [label=\textbf{\arabic*.}, leftmargin=0em]

      \item Let $X$ be a surface in $\P^n$, $n \geq 3$, defined as the complete intersection of hypersurfaces of degree $d_1, \dots, d_{n - 2}$, with each $d_i \geq 2$.
            Show that for all but finitely many choices of $(n, d_1, \dots, d_{n - 2})$, the surface $X$ is of general type.
            List the exceptional cases, and where they fit into the classification picture.

            \begin{proof}
                  We will use (II, Ex. 8.4e), which states the canonical sheaf $\omega_X$ is isomorphic to $\fO_X(\sum_{i = 1}^{n - 2} d_i - n - 1)$, and assume $d_i$ is in descending order.
                  Since $d_i \geq 2$, for every $(n, d_1, \dots, d_{n - 2})$ with $n \geq 6$, the inequality $$\sum^{n - 2}_{i = 1} d_i - n - 1 > 0$$ always holds, which implies $\omega_X$ is very ample, hence $X$ is of general type.
                  If $n = 3$ and $d_1 \geq 5$, then $X$ is of general type.
                  If $n = 4$ and either $d_1 \geq 4$ or $d_2 \geq 3$, then $X$ is of general type.
                  If $n = 5$ and $d_1 \geq 3$, then $X$ is of general type.
                  Thus, the exceptional cases are:

                  \fbox{$n = 3, d_1 = 2$} $X$ is the nonsingular quadric surface, isomorphic to $\P^1 \times \P^1$, and it is also the rational ruled surface with invariant $e = 0$.

                  \fbox{$n = 3, d_1 = 3$} $X$ is a cubic surface, which is isomorphic to $\P^2$ with 6 points in general position blown up, and it is also called a Del Pezzo surface of degree 3.

                  \fbox{$n = 3, d_1 = 4$} In this case, $\omega_X \cong \fO_X$, and we have a short exact sequence
                  \[ \begin{tikzcd}
                              0 \arrow[r] & \fO_{\P^3}(-4) \arrow[r] & \fO_{\P^3} \arrow[r] & \fO_{X} \arrow[r] & 0
                        \end{tikzcd} ,\]
                  and the long exact sequence in cohomology gives
                  \[ \begin{tikzcd}
                              \cH^1(\P^3, \fO_{\P^3}) \arrow[r] & \cH^1(X, \fO_X) \arrow[r] & \cH^2(\P^3, \fO_{\P^3}(-4))
                        \end{tikzcd}, \]
                  so $q = h^1(X, \fO_X) = 0$.
                  Hence, $X$ is a K3 surface.

                  \fbox{$n = 4, d_1 = d_2 = 2$} $X$ is a Del Pezzo surface of degree 4 with $\omega_X \cong \fO_X(-1)$, which is isomorphic to $\P^2$ with 5 points blown up (Ex. 4.13).

                  \fbox{$n = 4, d_1 = 3, d_2 = 2$} In this case, $\omega_X \cong \fO_X$.
                  Let $\fI$ be the ideal sheaf of $X$ in $\P^4$.
                  The long exact sequence in cohomology gives the isomorphism
                  \begin{equation*}
                        \cH^1(X, \fO_X) \cong \cH^2(\P^4, \fI).
                  \end{equation*}
                  By (III, \SS 7), we can use the Koszul resolution of $\fI$ to compute its cohomology
                  \[ \begin{tikzcd}
                              0 \arrow[r] & \fO_{\P^4}(-5) \arrow[r] & \fO_{\P^5}(-3) \oplus \fO_{\P^4}(-2) \arrow[r] & \fI \arrow[r] & 0
                        \end{tikzcd} ,\]
                  which readily shows $\cH^2(\P^4, \fI) = 0$.
                  Hence, $X$ is a K3 surface.

                  \fbox{$n = 5, d_1 = d_2 = d_3 = 2$} In this case, $\omega_X \cong \fO_X$.
                  Let $\fI$ be the ideal sheaf of $X$ in $\P^5$.
                  The long exact sequence in cohomology gives the isomorphism
                  \begin{equation*}
                        \cH^1(X, \fO_X) \cong \cH^2(\P^5, \fI).
                  \end{equation*}
                  By (III, \SS 7), we can use the Koszul resolution of $\fI$ to compute its cohomology
                  \[ \begin{tikzcd}
                              0 \arrow[r] & \fO_{\P^5}(-6) \arrow[r] & \fO_{\P^5}(-4)^{\oplus 3} \arrow[r] & \fO_{\P^5}(-2)^{\oplus 3} \arrow[r] & \fI \arrow[r] & 0
                        \end{tikzcd} ,\]
                  By (III, \SS 7), we can use the Koszul resolution of $\fI$ to compute its cohomology
                  which readily shows $\cH^2(\P^5, \fI) = 0$.
                  Hence, $X$ is a K3 surface.
            \end{proof}

      \item Prove the following theorem of Chern and Griffiths.
            Let $X$ be a nonsingular surface of degree $d$ in $\P_\C^{n + 1}$, which is not contained in any hyperplane.
            If $d < 2n$, then $p_g(X) = 0$.
            If $d = 2n$, then either $p_g(X) = 0$, or $p_g(X) = 1$ and $X$ is a K3 surface.

            \begin{proof}
                  
            \end{proof}
\end{enumerate}

\end{document}
